\chapter{Discussion}\label{ch:discussion}

This chapter maps the results back to the research questions and discusses how far the conclusions generalise beyond the DYNAMOS case study. The emphasis is on what the measurements mean for policy-aware, dynamically composed data exchange systems, rather than on selecting a single transport in isolation.

\section{RQ1: Transport Suitability}\label{sec:discussion-rq1}

RQ1 asks which streaming communication technologies, frameworks, and interaction patterns are suitable for large-scale and continuous data exchange in microservice-based systems. The answer is boundary-specific. For local service-to-service transfer, direct gRPC mappings are strongest. Client-streaming gives the highest raw clean throughput, but batched unary is close enough to change the interpretation of the unary result. The weak per-message unary column is not a verdict against request-response communication. It is a verdict against one-request-per-message transfer. Once multiple records are carried per request, unary becomes a strong bulk-transfer mechanism.

Brokered transports remain useful design evidence even though they are slower in the local clean matrix. RabbitMQ Streams is not validated here as a local latency improvement over chunked unary. Its value is architectural: retained records, offset-based replay, and decoupled consumption may matter more in a distributed deployment than on a single node. There is also a compliance-specific argument that the clean matrix cannot show: a retained stream is a replayable record of exactly which chunks crossed an agent boundary and in which order. In a policy-aware system, the transfer log itself can serve as an audit artifact, allowing a data steward to reconstruct what was delivered under an agreement after the fact. Direct transports provide nothing equivalent without additional logging infrastructure. NATS and Kafka are informative comparison points, but both introduce recovery, ordering, configuration, and platform shifts that make them lower-priority choices for the current DYNAMOS implementation.

The resulting RQ1 answer is therefore not a universal ranking. The best raw transport is client-streaming gRPC. The best request-shaped bulk-transfer pattern is batched unary. The best implementation hypothesis for the case study is chunked transfer at existing boundaries, with gRPC streaming as an in-pod optimisation and RabbitMQ Streams as a distributed-boundary candidate rather than as the source of the local speedup.

\section{RQ2: Integration with Policy-Aware Orchestration}\label{sec:discussion-rq2}

RQ2 asks how streaming communication can be integrated while preserving dynamic service composition, policy-driven orchestration, and isolation. The DYNAMOS implementation answers this by changing the bulk data path after policy approval while leaving the control path intact. The implementation introduces an explicit chunk model at the SQL query service, preserves those chunks through aggregate and algorithm services, carries them through the agent boundary, and exposes partial result events to the client.

This design preserves the important control properties of the case-study system. Requests are still evaluated before execution, generated jobs remain the execution unit, and provider selection remains policy-driven. The streaming contribution is therefore deliberately scoped: it keeps the pre-decision enforcement model of the original system and does not add continuous, usage-control-style enforcement~\cite{park2004ucon}. As \Cref{sec:policy-implications} argues, however, the chunk model changes what enforcement can do: every chunk boundary is a potential enforcement point, revocation already works at job granularity, and an interrupted transfer leaves an auditable delivered prefix rather than an undefined state. The mechanism for chunk-boundary re-validation is in place even though the policy wiring is future work.

Two alternative integration designs were considered and rejected for this thesis. A sidecar-as-stream-proxy design, in which generated services stay unary and sidecars translate between unary calls and streams in the style of Dapr~\cite{figueira2024selfadaptive}, would concentrate streaming complexity in infrastructure code, but it hides chunk boundaries from the services that must preserve them and would have required deeper sidecar changes than the evaluated variants. A claim-check design, in which chunks are written to shared storage and only references travel through the approved path, is the standard enterprise answer for very large payloads, but it moves the actual data outside the policy-mediated communication path and shifts the trust boundary to storage access control, which conflicts with the DYNAMOS premise that the mediated path is the enforcement surface.

The best architectural choice is again boundary-specific. Chunked unary is the strongest current local default because it gives chunking without changing the DYNAMOS execution model very much. Intra-job gRPC streaming is useful where both endpoints are already in the generated pod and a direct stream is acceptable. RabbitMQ Streams is more relevant at the inter-agent boundary, especially for future distributed deployments where retained streams, replay, and independent consumers may matter more than local latency.

\section{RQ3: Case-Study Impact}\label{sec:discussion-rq3}

RQ3 asks what impact the implemented streaming data paths have on scalability, performance, and robustness compared to traditional request-response transfer. The DYNAMOS-specific benchmark changes the interpretation of the isolated transport study. RQ1 showed that client-streaming gRPC is the strongest raw transport in the controlled pipeline. The DYNAMOS evaluation shows that integration fit matters as much as raw transport speed.

The central hypothesis survives validation in a more precise form: chunking is the main source of the improvement, not the name of the transport alone. All chunk-aware DYNAMOS data paths improve local bulk-retrieval scalability compared with classic unary. They reduce completion time, expose earlier partial results, and lower peak memory usage for 50k, 100k, and 250k row result sets. The best current local default is chunked unary. Intra-job gRPC streaming is a close alternative for pod-local service chains. RabbitMQ Streams should be kept as a distributed-boundary candidate, but its advantage must be validated under cross-site conditions rather than inferred from the local latency result.

\section{Comparison with Existing Literature}\label{sec:discussion-literature}

The results agree with prior work at the level of communication tradeoffs, while the absolute numbers are not directly comparable. Coviello et al.~\cite{coviello2022dataxc} show in DataXc that flexible broker-like communication and direct microservice communication involve different latency, throughput, and network-cost tradeoffs. The thesis finds the same pattern in a policy-aware data exchange setting: direct gRPC transfer is strongest locally, while brokered streams remain relevant for decoupling and replay.

Roohitavaf et al.~\cite{roohitavaf2019logplayer} show that asynchronous gRPC streaming can outperform a Kafka-based design for ordered log delivery under their workload. This supports the thesis observation that direct streaming is strong when endpoints and ordering requirements are tightly scoped. Dobbelaere and Sheykh Esmaili~\cite{dobbelaere2017kafkarabbitmq} and Ibrahim et al.~\cite{ibrahim2026brokers} both show that broker rankings depend on configuration, workload, and reliability requirements. This is consistent with the thesis finding that RabbitMQ, NATS, and Kafka should not be presented as universally ordered by one local matrix.

Marcu and Bouvry~\cite{marcu2022pushpull} emphasise that push/pull choices and consumer-controlled flow affect stream-processing throughput and latency. The backpressure and recovery scenarios in this thesis point in the same direction: pacing and overload behaviour are as important as clean throughput. The thesis-specific addition is the batched unary result. It shows that a request-shaped design can become competitive when it preserves chunks and amortises per-message overhead, which is especially relevant for systems that already rely on policy-approved request orchestration.

\section{Threats to Validity}\label{sec:threats-validity}

\subsection{Transport benchmark threats}

Four limits are important when interpreting the transport benchmark. First, all primary RQ1 evidence comes from one local single-node Kubernetes environment. It is suitable for isolating transport overhead, but it does not capture cross-node latency, broker placement, replication, or geographically distributed deployment. Second, the benchmark uses synthetic payloads and fixed request shapes. This makes the transports comparable, but it also means that the RQ1 results should be read as mechanism-level evidence rather than as a complete data exchange workload evaluation. Third, the tables report arithmetic means over five repetitions. Kubernetes scheduling, broker flushing, and container resource sharing can still introduce run-to-run variation. Small differences between brokered transports should therefore not be treated as decisive. Fourth, literature comparison is directional rather than numeric: prior studies use different workloads, broker modes, deployment topologies, and persistence settings.

\subsection{DYNAMOS benchmark threats}

The RQ3 evidence has five main limits. First, the benchmark runs on a local Kubernetes deployment, so it does not capture cross-site latency, bandwidth variation, or independent administrative domains. Second, the additional scenario checks cover two alternative request shapes, not the full space of DYNAMOS archetypes, query forms, and provider combinations. They support compatibility, but they do not replace a broader workload campaign. Third, resource usage is sampled once per second, so peak CPU and memory values can miss short spikes. Fourth, the fixed 5{,}000-row chunk size is justified experimentally but not optimised across all workloads. Fifth, the measured request window includes policy handling, orchestration, generated job startup, and response serialisation. That is fair for comparing user-visible DYNAMOS behaviour, but it compresses the apparent transfer-only speedup because fixed control-path costs remain present in every implementation.
